% ============================================================
%  Setmana 34 — Unitat 11 · Integrals
%  Exercicis: 52, 54
% ============================================================

\setmana{34}{Unitat 11 \textperiodcentered{} Integrals}{9 de maig de 2027}%
        {52 i 54}

\begin{errada}{Exercici 52 g)}
  \ondiu{$\displaystyle\int (x-1)^3\,dx=\dfrac{(x-1)^3}{4}+k$}
  \hadedir{$\displaystyle\int (x-1)^3\,dx=\dfrac{(x-1)^4}{4}+k$}
  \motiu{En integrar una potència, l'exponent augmenta una unitat.}
\end{errada}

\begin{errada}{Exercici 54 a)}
  \ondiu{$\displaystyle\int\dfrac{4}{x-2}\,dx=4\ln|x+2|+k$}
  \hadedir{$\displaystyle\int\dfrac{4}{x-2}\,dx=4\ln|x-2|+k$}
  \motiu{El denominador és $x-2$.}
\end{errada}

\begin{errada}{Exercici 54 d)}
  \ondiu{$\displaystyle\int\dfrac{1}{(x+4)^2}\,dx=-\dfrac{1}{x-4}+k$}
  \hadedir{$\displaystyle\int\dfrac{1}{(x+4)^2}\,dx=-\dfrac{1}{x+4}+k$}
  \motiu{El denominador és $x+4$.}
\end{errada}

\begin{errada}{Exercici 54 f)}
  \ondiu{$\displaystyle\int\left(\dfrac{2}{(x+3)^3}+\dfrac{5}{(x-3)^2}\right)dx
    =\dfrac{-1}{(x+3)^2}-\dfrac{5}{x-3}+k$}
  \hadedir{$\displaystyle\int\left(\dfrac{2}{(x+3)^3}-\dfrac{5}{(x-3)^2}\right)dx
    =\dfrac{-1}{(x+3)^2}+\dfrac{5}{x-3}+k$}
  \motiu{A l'enunciat, el segon terme porta signe menys. La solució el copia
    amb signe més, i el resultat que dona és el d'aquesta altra integral.}
\end{errada}
